\documentclass[11pt,a4paper]{article}
\usepackage[T1]{fontenc}
\usepackage{lmodern,amsmath,amssymb,booktabs,microtype}
\usepackage[margin=24mm]{geometry}
\usepackage[colorlinks=true,urlcolor=blue,linkcolor=blue]{hyperref}
\newcommand{\Js}{\mathcal J_s}
\newcommand{\Aop}{\mathcal A}
\newcommand{\HH}{\mathcal H}
\newcommand{\MM}{\mathcal M^2}
\DeclareMathOperator{\sech}{sech}
\title{The ordering-odd inner product\\A relative Casimir and its translation test}
\author{Alan Day \and Chengheng\thanks{AI research assistant (OpenAI).}}
\date{Local research note --- 6 November 2026}
\begin{document}\maketitle
\begin{abstract}
The known leading commutator four-point function determines a positive
two-particle quadratic form: the difference between the total and the
two one-particle Casimirs. We derive this directly from coherent-state
matrix elements, without assuming the conjectured all-degree finite-$q$
ordering overlap. For the quadratic and cubic singlets the relative
eigenvalues are $k(k+9)+12$. We construct self-adjoint light-front
translations in this weighted space and then test the stronger condition
that they preserve the microscopic product vertex. An explicit norm
calculation distinguishes these two requirements.
\end{abstract}
\section{The leading inner product}
Use the annealed normalized trace of complex Gaussian $p$-body DSSYK,
standard CAR including overlapping index sets, and
$\beta=\mu_{\rm chem}=0$. First take $N,p\to\infty$ at
$\lambda=p^2/N>0$, then small $h=4\lambda$, with
$q=e^{-h}$, $u=2\Js t$ and $\hbar=1$. The normalized derivative
singlets of the preceding notes have labels $d=2,e=3$ \cite{local}.

The free limiting product representation is
$\HH_0=D_d^+\otimes D_e^+$. Write its generators as $E_j,B_j,\Pi_j$,
with $C_j=E_j^2-B_j^2-\Pi_j^2=d_j(d_j-1)$ and primitive total
generators. The leading norm of the rescaled commutator is
\begin{equation}
 \boxed{\langle f,g\rangle_{\rm odd}=\langle f,\Aop g\rangle_0,
 \qquad \Aop=C_{\rm tot}-C_1-C_2
 =2(E_1E_2-B_1B_2-\Pi_1\Pi_2).}\label{metric}
\end{equation}
For $d=2,e=3$, the relative-channel spectrum is
\begin{equation}
 \boxed{\Aop_k=k(k+9)+12,\qquad k=0,1,\ldots.}\label{spectrum}
\end{equation}
This is the entire positive quadratic form of the leading theory, not
a finite-$q$ all-$k$ overlap formula. It identifies the norm required
by the extra ordering channel before a spatial dynamics is chosen.

Lin--Stanford already give the leading commutator kernel and its block
weights \cite{LS}. We establish their operator realization for the
present singlets and use it to examine translations. No priority claim
for the commutator sector or the discrete-series decomposition is made.

\newpage
\section{Direct identification from the known kernel}
Realize $D_d^+$ on $L^2(\mathbb R_+,dp)$ with
\begin{gather}
 Q=p,\qquad B=i(p\partial_p+\tfrac12),\qquad
 S=-p\partial_p^2-\partial_p+(d-\tfrac12)^2/p,\nonumber\\
 E=(Q+S)/2,\qquad \Pi=(Q-S)/2,\qquad
 \psi_d(p)=\frac{2^d}{\sqrt{\Gamma(2d)}}p^{d-1/2}e^{-p}.
\end{gather}
The states $|u\rangle=e^{-2iuB}\psi_d$ have real wavefunctions
$e^u\psi_d(e^{2u}p)$ and overlap
$g_d(u-u')=\sech^{2d}(u-u')$.
Elementary gamma integrals give
\begin{align}
 \frac{\langle u'|E|u\rangle}{g_d(u-u')}
 &=d\frac{\cosh(u+u')}{\cosh(u-u')},&
 \frac{\langle u'|B|u\rangle}{g_d(u-u')}
 &=-id\tanh(u-u'),\nonumber\\
 \frac{\langle u'|\Pi|u\rangle}{g_d(u-u')}
 &=-d\frac{\sinh(u+u')}{\cosh(u-u')}.
\end{align}
Taking tensor products and substituting in \eqref{metric} yields
\begin{align}
 \langle u',v'|\Aop|u,v\rangle
 &=2de\,\sech^{2d+1}(u-u')\sech^{2e+1}(v-v')\nonumber\\
 &\quad\times\bigl[\cosh(u+u'-v-v')
             +\sinh(u-u')\sinh(v-v')\bigr].\label{kernel}
\end{align}
This reproduces the covariance derived in the preceding note from
Lin--Stanford Eqs.~(170), (205) and (388):
\[
 B_h(u,v)=-i[V(u),W(v)]/\sqrt h,\qquad
 \lim_{h\to0}\tau\bigl(B_h(u',v')B_h(u,v)\bigr)
 =\langle u',v'|\Aop|u,v\rangle.
\]
The equality includes arbitrary finite linear combinations of time
insertions. It uses the sourced leading fixed-time kernel, and does
not require differentiating a numerically fitted function.

\paragraph{Closure and completeness.}
The tensor decomposition \cite{KV} reads
$D_d^+\otimes D_e^+=\bigoplus_{k\ge0}D_{d+e+k}^+$.
Thus $\Aop_k=k(k+2d+2e-1)+2de\ge2de>0$.
Its spectral realization defines a positive self-adjoint operator.
Finite tensor descendant sums are a form core. Derivatives at zero of
the boost orbits span these descendants: $B^n\psi_d$ is triangular in
the ladder basis with a nonzero top coefficient. The analytic orbit
vectors admit these difference quotients also in the form norm.
Consequently their span is form-dense, and the kernel identifies the
closed leading form uniquely. This is not a claim of convergence
uniform in channel number or in times growing with $h^{-1}$.

\newpage
\section{An explicit relative differential operator}
Use $P=p_1+p_2>0$, $x=p_1/P\in(0,1)$. After extracting the square root
of the beta weight, the product Hilbert space is
\begin{equation}
 \HH_0=L^2(dP\,w(x)dx),\qquad
 w(x)=\frac{x^{2d-1}(1-x)^{2e-1}}{\mathrm B(2d,2e)}.
\end{equation}
The Jacobi--Laguerre tensor map \cite{KV,local} gives
\begin{equation}
 \Aop=-x(1-x)\partial_x^2-[2d-2(d+e)x]\partial_x+2de.\label{diff}
\end{equation}
Its closed quadratic form, initially on finite Jacobi sums, is
\begin{equation}
 \|f\|_{\rm odd}^2=\int dP\int_0^1w(x)
 \left[x(1-x)|\partial_x f|^2+2de|f|^2\right]dx.\label{form}
\end{equation}
The normalized relative eigenfunctions are Jacobi polynomials
$p_k(x)\propto P_k^{(2e-1,2d-1)}(2x-1)$. At $d=2,e=3$,
\[
 w(x)=504x^3(1-x)^5,\quad
 \Aop=-x(1-x)\partial_x^2-(4-10x)\partial_x+12,
\]
and its first six eigenvalues are $12,22,34,48,64,82$.
The equal-time normalized commutator of the last note is the $k=0$
seed, divided by $\sqrt{12}$ in this form norm.

Define $\HH_{\rm odd}=D(\Aop^{1/2})$, completed in \eqref{form}.
Since $\Aop\ge12$, it is already complete in that norm, and
\begin{equation}
 U_{\Aop}:\HH_{\rm odd}\to\HH_0,\qquad U_{\Aop}f=\Aop^{1/2}f
\end{equation}
is unitary and onto. Its inverse is bounded on $\HH_0$.
The total generators $E,B,Q$ commute strongly with $\Aop$ in the
direct-sum representation. Their actions therefore carry over unchanged
to the weighted space; specifically
$Q=P$ and $B=i(P\partial_P+\tfrac12)$.

\paragraph{The product vertex in normalized coordinates.}
Writing an ordered product as half the anticommutator plus half the
commutator, the leading map into the normalized odd copy is
\begin{equation}
 \mathcal V_h f=\frac{i\sqrt h}{2}\Aop^{1/2}f.\label{vertex}
\end{equation}
On $k=0$ it is $i\sqrt{3h}=i\sqrt{12\lambda}$, reproducing the
previous three-point overlap. The amplitude grows with relative $k$;
the fixed-channel expansion is not uniform for $hk^2$ of order one.
The closed form specifies which wave packets have finite leading
production norm: exactly $f\in D(\Aop^{1/2})$.

\newpage
\section{A norm-preserving translation construction}
For positive constituent masses $m_1,m_2$ the standard free light-front
mass profile \cite{Pauli} is
\begin{equation}
 M(x)=\frac{m_1^2}{x}+\frac{m_2^2}{1-x},\qquad
 P^+=\mu P,\qquad P^-_0=\frac{M(x)}{\mu P},\qquad \mu>0.
\end{equation}
Here $P,x$ are dimensionless, whereas $m_j,\mu$ have energy units.
Multiplication by $M$ is positive self-adjoint in $\HH_0$, with spectrum
$[(m_1+m_2)^2,\infty)$. It cannot simply be reused in the odd metric:
\begin{align}
 [\Aop,M]f={}&-2x(1-x)M'f'\nonumber\\
 &-\bigl[x(1-x)M''+(4-10x)M'\bigr]f.\label{defect}
\end{align}
This is nonzero for the free profile. The coefficient of $f'$ shows
that a smooth multiplication profile commuting with $\Aop$ must be constant.

There is an explicit self-adjoint construction in the weighted space:
\begin{equation}
 \boxed{\MM_{\rm odd}=\Aop^{-1/2}M\Aop^{1/2},\qquad
        P^-_{\rm odd}=\frac{\MM_{\rm odd}}{\mu P}.}\label{transport}
\end{equation}
Its domain is the inverse image under $U_{\Aop}$ of the multiplication
domain of $M$ (and of $M/(\mu P)$ for $P^-$). Unitary transport proves
self-adjointness, positivity, and the same mass spectrum. It also gives
strongly commuting $P^+,P^-$ and
\[
 [B,P^+]=iP^+,\qquad [B,P^-]=-iP^-,\qquad [P^+,P^-]=0
\]
on transported smooth cores. This solves the Hilbert-space consistency
of a candidate free translation; it does not yet match the product vertex.

\paragraph{An exact diagnostic.}
Set $m_1=m_2=1$ in a fixed-$P$ fiber, and use
$p_0=1$, $p_1=(10x-4)/\sqrt{24/11}$.
Then $\langle p_0,Mp_0\rangle_0=24/5$ and
$\langle p_0,Mp_1\rangle_0=-\sqrt{66}/10$.
For $f=(p_0+ip_1)/\sqrt2$, direct integration gives
\begin{equation}
 \boxed{\|e^{-isM}f\|_{\rm odd}^2=17+\sqrt{66}\,s+181s^2.}\label{norm}
\end{equation}
Thus the undressed multiplication phase genuinely fails to preserve
the physical odd norm; this is not a finite matrix truncation effect.
The transported evolution \eqref{transport} preserves it exactly.

\newpage
\section{The stronger product-compatibility condition}
Let $T=\Aop^{1/2}$, so the microscopic leading vertex in normalized
coordinates is $i\sqrt h\,T/2$. Suppose it is to be a translation
intertwiner between an ordinary two-particle mass $M$ and some
self-adjoint odd-sector mass $N$:
\begin{equation}
 e^{-isN}Tf=Te^{-isM}f,\qquad f\in D(T),\label{intertwine}
\end{equation}
with domain invariance. Taking norms forces
$\|Te^{-isM}f\|_0=\|Tf\|_0$, i.e. invariance of the $\Aop$ form under
the even-sector evolution. Equation \eqref{norm} contradicts this
for the primitive free profile. No choice of self-adjoint $N$ on this
closed target sector can fix that particular intertwining requirement.

Equivalently, on a suitable common core the necessary condition is
$[\Aop,M]=0$. In normalized relative coordinates,
\begin{equation}
 [\Aop,M]_{01}=\sqrt{66},\qquad
 [T,M]_{01}=(\sqrt{12}-\sqrt{22})(-\sqrt{66}/10)\ne0.
\end{equation}
The transported odd mass on page~4 becomes simply $M$ after applying
$U_{\Aop}$. It therefore preserves norms but fails \eqref{intertwine}.
The two statements have different requirements and must not be conflated.

\paragraph{What can satisfy both conditions?}
A positive relative mass operator $F(\Aop)$ commutes with the vertex,
and yields kinematically consistent translations and intertwining on
both copies. Its masses $F(k(k+9)+12)$ remain arbitrary; choosing such
a function is not a derivation of meson dynamics. More generally the
assumptions behind \eqref{intertwine} may have to change: the microscopic
product insertion need not already be an on-shell transition map, and
physical translations may involve additional sectors or a different
spatial operator dictionary. The conditional incompatibility concerns
the specified free-pair identification, not the existence of a bulk dual.

\section{Checks of the actual claims}
Fifty exact symbolic identities verify the coherent Casimir kernel,
general differential commutator, Jacobi conjugation, six relative
eigenfunctions, their full orthogonality matrix, and the translation
diagnostic. The $6\times6$ free mass matrix is retained.
Five independent relative-coordinate integrals at 55 and 85 digits
reproduce the four-time kernel; the largest 85-digit absolute residual
is $5.15\,10^{-85}$. Four 70-digit integrals check \eqref{norm}, with
residual below $2.90\,10^{-70}$. These assess numerical precision;
the displayed identities also have analytic derivations.

\newpage
\appendix
\section{Finite-$h$ controls and the next physical input}
Independent finite-$q$ direct/cross Gram recursions construct primaries
through $k=4$ for $d=2,e=3$, at $h=.04,.02,.01$. For the parity-adjusted
overlap $\rho_k$, compare $2(1-\rho_k)/h$ with \eqref{spectrum}.
At $h=.01$ the results are
\begin{center}
\begin{tabular}{rrr}
\toprule
$k$&finite-$h$ odd norm divided by $h$&leading eigenvalue\\\midrule
0&11.64709328&12\\
1&20.83317294&22\\
2&31.26703668&34\\
3&42.67442779&48\\
4&54.77019259&64\\\bottomrule
\end{tabular}
\end{center}
Increasing precision from 85 to 115 digits changes the coefficients
relatively by at most $4.37\,10^{-79}$. Nevertheless the finite-$h$
deviation at $k=4$ is still $14.42\%$. Numerical precision does not
remove the finite-deformation correction.

Lin--Stanford Appendix F describes its general finite-$q$ overlap
formula, Eq.~(355), as inferred from small-degree calculations.
Our all-$k$ leading form is instead obtained from the independently
available leading commutator kernel, its coherent-state realization,
and the standard tensor Casimir. The low-degree finite-$q$ comparisons
are checks, not the proof of that limiting operator. We make no uniform
large-$k$ or finite-$h$ operator-norm convergence assertion.

\paragraph{Next research step.}
The leading ordering metric and its domain are now explicit. A free
translation can be made unitary, but its compatibility with the
microscopic multiplication data requires extra physics. The next step
is to return to the Susskind bulk gauge/meson construction and identify
which spatial insertions or interacting translation generators its
dictionary actually proposes. The tests are the norm \eqref{form},
the vertex \eqref{vertex}, and the conditional Ward relation
\eqref{intertwine}. We should derive or test that map, rather than choose
an arbitrary function of the Casimir and call it a mass spectrum.

\begin{thebibliography}{9}
\bibitem{LS} H. W. Lin and D. Stanford, \emph{A symmetry algebra in
double-scaled SYK}, \href{https://arxiv.org/abs/2307.15725}{arXiv:2307.15725},
Section 4.3, Eq.~(170), and Appendix F.1.
\bibitem{KV} H. T. Koelink and J. Van der Jeugt,
\emph{Convolutions for orthogonal polynomials from Lie and quantum algebra
representations}, \href{https://arxiv.org/abs/q-alg/9607010}{q-alg/9607010},
Section 3, Eq.~(3.2), Corollary 3.6 and Remark 3.7.
\bibitem{Pauli} H.-C. Pauli, \emph{A compendium of light-cone quantization},
\href{https://arxiv.org/abs/hep-ph/0103106}{hep-ph/0103106}, Eq.~(87).
\bibitem{local} A. Day and Chengheng, local notes \emph{Two matter},
\emph{Antisymmetric descendants}, and \emph{Commutator primary}
(November 2026). Evidence: \texttt{results/susskind-ordering-metric}.
\end{thebibliography}
\end{document}
